\section{The four residual sextic discriminants}\label{sec:four-sieves}

It remains to treat $D=87,348,432,675$, where $D=|\Delta|$.
Their degree-six class polynomials were computed by enclosing the product
over all six reduced forms. Apply Lemma~\ref{lem:confluent} to five powers
of a root $x$. After dividing by the lowest power, one exponent is zero.
In the basis $1,x,\ldots,x^5$ its column is $1$. Omitting the constant
coordinate therefore leaves a $5\times4$ matrix. At the prime $p$, a
residue with multiplicity $l$ among \emph{all five} exponents contributes
the confluent columns of orders $0$ through $l-1$, with the distinguished
column removed afterwards. A candidate survives exactly when every
four-by-four minor vanishes at every prime tested at the current stage.
This gives finite tests even when the integer exponents grow without bound
or two of them coincide modulo arbitrarily high prime powers.

The table records the complete enumeration. The period $m$ in each
pair is certified by reduction modulo $(H_{-D},p^2)$ as in
Theorem~\ref{thm:131}. Each step from $n$ to a multiple $N$ enumerates
all lifts of the preceding sorted quadruples. All four nonzero-position
residues may themselves be zero.
\begin{center}\footnotesize
\begin{tabular}{rrrrrr}
\toprule
$D$ & $n$ & $N$ & $(p,m)$ & Tested & Left\\
\midrule
87 & 1 & 102 & $(103,102)$ &4,780,230&13,125\\
 &102 &1,530 & $(1531,1530)$ &653,917,425&625\\
 &1,530&7,650& $(151,150)$ &390,625&915\\
 &18&486& $(487,486),(1459,486)$ &7,971,615&0\\
348&1&486& $(487,486),(1459,486)$ &2,353,329,126&0\\
432&1&36& $(109,36)$ &82,251&106\\
 &36&144& $(433,144)$ &26,212&0\\
675&1&40& $(241,40)$ &123,410&1\\
 &40&280& $(13,56)$ &2,401&0\\
\bottomrule
\end{tabular}
\end{center}
For $D=87$ the 915 survivors modulo $7650$ reduce to precisely 15 sorted
quadruples modulo $18=\gcd(7650,486)$. The last step for $87$ tests every
lift of those 15 possibilities to $486$. Any set of integer exponents would
give compatible residues on both sides, so the zero final count excludes
every such set. For the other discriminants the successive moduli already
divide one another. In all four cases any five powers are independent over
$\Q$.

Run \texttt{python3 checks/verify\_five\_sieves.py} from the archive root.
The program reads the certified integer class polynomial coefficients,
constructs every finite-field column, compiles the C++17 enumerator, and
checks each count. The enumerator streams the candidates and stores only
the surviving quadruples. Thus even the stage testing more than two
billion candidates needs room for only its residue tables and survivors.
The full direct replay took 33 seconds including compilation on the
verification machine. Its largest list had 13,125 quadruples, and
the largest coefficient table had $1530\cdot5\cdot5$ field entries.

A second implementation, \texttt{checks/fast\_sieve/verify.py}, checks $D=87,348$ by eliminating two prefix columns and grouping suffix columns by their projective classes in the quotient. For $D=87$, the $196{,}587$ simultaneous failures modulo 486, the $13{,}125$ failures modulo 102 and the $19{,}669$ failures modulo 150 restrict the lifts to 979,524 candidates at modulus 1530. The rank test at $1531$ eliminates all of them. For $D=348$ the independent calculation already finds zero modulo 486.

Combining these four exclusions with the preceding 67, the unbounded-exponent quintic test and the global cutoff proves the reverse implication in Theorem~\ref{thm:four}. For $h\le4$, the degree of $H_\Delta$ proves the forward implication.
